\documentclass[conference]{IEEEtran}
\IEEEoverridecommandlockouts

\usepackage{cite}
\usepackage{amsmath,amssymb,amsfonts}
\usepackage{graphicx}
\usepackage{textcomp}
\usepackage{xcolor}
\usepackage{url}
\usepackage{booktabs}
\usepackage{array}

\begin{document}

\title{Measuring the Art.\,52 Gap: A GitHub-Scale Empirical Study of EU AI Act
Transparency Compliance in AI-Assisted Software Development}

\author{%
  \IEEEauthorblockN{Vinita Silaparasetty}
  \IEEEauthorblockA{%
    \textit{Aevoxis Solutions}\\
    info@aevoxis.de}
  \thanks{Manuscript submitted 2 July 2026.
    Aggregate experimental results (GitHub API counts, header survival
    outcomes, PyPI audit totals) are available at
    \url{https://github.com/VinitaSilaparasetty/pr-automation-agent}.
    No individual repository names, usernames, or personal identifiers
    are included in the dataset.}
}

\maketitle

\begin{abstract}
The EU Artificial Intelligence Act (EU 2024/1689) requires under
Article~52 that AI-generated content be labelled to enable human
oversight.  We present a GitHub-scale empirical study of whether this
obligation is being met across publicly available Python repositories.
Three experiments on freely accessible public data address: (RQ1) the
fraction of repositories actively configured for GitHub Copilot that
implement any Art.\,52 label; (RQ2) whether such labels persist through
standard Python code transformation pipelines; and (RQ3) whether any
compliance signal survives into PyPI package distributions.  Across
43,816 Copilot-configured repositories and 18,210 Python files from
25~major PyPI packages, spontaneous compliance is near zero and
provenance is entirely absent from the distributed supply chain.
Comment-based labels survive standard formatters but are destroyed by
bytecode compilation.  We identify two under-examined regulatory blind
spots---provenance escape in the software supply chain and output-context
risk elevation in risk classification---situate Art.\,52 within the
EU's three-instrument AI regulatory package, compare it with China's
Deep Synthesis Regulations, and derive four concrete policy amendments.
\end{abstract}

\begin{IEEEkeywords}
EU AI Act, Article~52, AI transparency, empirical compliance, GitHub
mining, software supply chain, AI Liability Directive, code provenance,
open source
\end{IEEEkeywords}

% =========================================================================
\section{Introduction}
% =========================================================================

The European Union Artificial Intelligence Act (EU~2024/1689), which
entered into force on 1~August 2024, is the world's first binding
horizontal AI regulation.  Its Article~52 establishes that providers
and deployers of AI systems that generate synthetic content---including,
per Recital~132, software source code---must ensure outputs are labelled
as AI-generated ``in a clear and distinguishable manner.''  The stated
objective is to enable meaningful human oversight before AI-generated
content is acted upon.

The obligation is the right instinct.  GitHub's State of the
Octoverse~2023 survey found that 92\% of enterprise developers now use
AI coding assistants at least some of the time~\cite{github2023}, and
independently replicated research has shown that developers using AI
coding assistants are significantly more likely to introduce security
vulnerabilities---and to do so with higher confidence---than developers
writing code manually~\cite{perry2023}.  The case for mandated
transparency is sound.  The question this paper asks is whether
Art.\,52 actually produces that transparency in practice.

Our answer, supported by three experiments conducted entirely on
publicly available, freely accessible data, is: it does not.  Across
43,816 public GitHub repositories that demonstrably use GitHub Copilot,
and across 18,210 Python files from 25 major PyPI package distributions,
the empirical Art.\,52 compliance rate is effectively zero.  The label
is not being written, and where it is written, it is invisible at the
distribution layer that most end-users encounter.  We argue that this
outcome is not a failure of individual compliance effort but a
predictable consequence of five structural deficiencies in the
regulatory text---deficiencies that the existing empirical literature on
formal versus substantive compliance in regulatory settings would have
predicted, and that targeted amendments can remedy.

The paper makes four contributions:
\begin{enumerate}
\item A GitHub-scale measurement of Art.\,52 compliance rates across
      Copilot-configured public repositories, the first such
      large-scale measurement to our knowledge.
\item An empirical header survival analysis across nine common Python
      code transformation tools and one compilation step.
\item A provenance loss measurement across 25 major PyPI package
      distributions.
\item A structured regulatory analysis situating Art.\,52 within the
      EU's three-instrument AI regulatory package and against China's
      technically more prescriptive Deep Synthesis Regulations.
\end{enumerate}

The remainder of the paper is organised as follows.
Section~\ref{sec:background} provides regulatory background and situates
the work in the literature.  Section~\ref{sec:rq} states research
questions and hypotheses.  Section~\ref{sec:methodology} presents
experimental methodology.  Section~\ref{sec:results} reports results.
Section~\ref{sec:blindspots} analyses two under-examined regulatory
blind spots.  Section~\ref{sec:doctrinal} develops the doctrinal
comparison.  Section~\ref{sec:recommendations} provides policy
recommendations.  Section~\ref{sec:threats} discusses threats to
validity.  Section~\ref{sec:conclusion} concludes.

% =========================================================================
\section{Background and Related Work}
\label{sec:background}
% =========================================================================

\subsection{EU AI Act Article~52 Transparency Obligations}

The EU AI Act stratifies AI systems by risk.  General Purpose AI (GPAI)
models---the category that covers GitHub Copilot and similar code
generation tools---are governed primarily by Articles~53--56, which
impose obligations on GPAI \emph{providers}.  Article~52 governs
\emph{deployers} of AI systems that interact with natural persons or
generate synthetic content.  For AI-assisted software development,
Article~52 applies to any organisation whose developers use GPAI coding
tools to generate Python (or any other language) files that are then
merged into production repositories.

Article~52(1) states that users must be informed, ``in a clear and
distinguishable manner,'' that content they are interacting with or
receiving is AI-generated, ``at the latest at the time of the first
interaction or exposure.''  Recital~132 makes explicit that this applies
to synthetic text outputs, which the Commission has since confirmed
includes generated source code.

Crucially, Article~52 specifies neither a technical format for the label
nor a retention period, schema, or enforcement mechanism for verifying
that labelling has occurred.  The Commission's AI Office, established
under Article~64 of the Act, has published guidance acknowledging that
technical standards for AI content marking are under development.  No
normative standard has been published as of the submission date of this
paper.

\subsection{Formal and Substantive Compliance: A Regulatory Theory Framework}

The distinction between formal compliance (the rule is followed in
letter) and substantive compliance (the rule's purpose is achieved) is
foundational in the sociology of law and organisational studies.
Edelman's seminal work on civil rights law in organisations demonstrated
that firms create formal structures---diversity offices, written
policies---that satisfy the letter of anti-discrimination law without
substantially changing hiring outcomes~\cite{edelman1992}.  The same
dynamic has been observed in GDPR cookie-consent
implementations~\cite{machuletz2020} and in HIPAA audit-log requirements
in healthcare IT systems.

Our paper applies this framework to Art.\,52: we ask whether the
compliance label achieves its stated purpose (meaningful human oversight
of AI-generated code) or merely satisfies the formal requirement (a
comment line in a file).  The empirical results speak directly to this
question.

\subsection{Developer Behaviour with AI-Assisted Tools}

Perry, Srivastava, Kumar, and Boneh~\cite{perry2023} conducted a
controlled experiment in which 47~participants---with varying programming
experience---were given security-sensitive coding tasks with and without
access to OpenAI Codex.  The Copilot-assisted group produced code with
significantly more security vulnerabilities ($p < 0.05$) and exhibited
higher confidence in their incorrect solutions.  The study did not
examine regulatory compliance infrastructure, but its implication for
Art.\,52 is direct: the productivity benefit of AI coding tools---faster
code acceptance---works against the substantive oversight that Art.\,52
attempts to mandate.

Chen et al.'s evaluation of Codex on the HumanEval benchmark~\cite{chen2021}
established that AI-generated code has correctness failures on
approximately 37\% of standard programming tasks before human review.
This establishes the \emph{need} for the kind of human oversight
Art.\,52 envisions; Perry et al.~\cite{perry2023} establishes that
developers using those same tools are less likely to provide it.  We
term this the Copilot Paradox and return to it in
Section~\ref{sec:blindspots}.

\subsection{Code Provenance and Software Supply Chain Security}

The SBOM (Software Bill of Materials) movement, codified in the SPDX
standard~\cite{spdx23} and mandated for US federal software procurement
by Executive Order~14028 (2021), provides an instructive analogy.  SBOM
requires that deployed software carry a machine-readable manifest
identifying all components and their origins.  AI-generated content
provenance is a strictly analogous problem: rather than identifying
component \emph{packages}, it requires identifying component
\emph{generation events}.  SPDX~3.0 introduced an experimental AI field
extension but has not been adopted for AI-generated code identification
at scale.

The C2PA (Coalition for Content Provenance and Authenticity)
specification~\cite{c2pa21} addresses media provenance using
cryptographically signed manifests that cannot be stripped without
breaking the artefact's integrity hash.  No equivalent standard exists
for source code.  The EU AI Act neither references nor mandates either
approach.

\subsection{The EU's Broader Regulatory Package}

The EU AI Act does not stand alone.  Two companion instruments interact
directly with Art.\,52.

The \textbf{AI Liability Directive (AILD, COM(2022)~496~final)}~\cite{aild2022}
proposes a disclosure-of-evidence mechanism (Art.\,3): where a plaintiff
demonstrates that an AI system was involved in causing damage and that
the defendant refused to disclose relevant evidence, the court may
presume the AI's causal role.  For AI-generated code, this creates a
material litigation risk: if a developer cannot produce an audit log
showing which files were AI-generated and who reviewed them, they lose
the evidentiary protection of the AILD's standard of care framework.

The \textbf{revised Product Liability Directive~((EU)~2024/2853)}~\cite{rpld2024}
extends product liability to software, closing the longstanding exclusion
of intangible goods.  AI-generated code that causes harm to a person or
property may now trigger producer liability for the entity that deployed
it.

Together with Art.\,52, these three instruments form a coherent (if
incompletely specified) package: Art.\,52 requires the label; the AILD
requires the audit trail to prove the label existed and review occurred;
the rPLD creates the liability exposure that gives both obligations their
enforcement context.  As Section~\ref{sec:doctrinal} will show, however,
the package is only coherent if Art.\,52 mandates the operational
elements (schema, retention, workflow controls) that would make the AILD
defence available.

% =========================================================================
\section{Research Questions and Hypotheses}
\label{sec:rq}
% =========================================================================

We address three research questions, each motivated by a distinct failure
mode in the Art.\,52 compliance chain.

\textbf{RQ1 (Voluntary Adoption):} What fraction of public GitHub
repositories that demonstrably use GitHub Copilot implement any form of
EU AI Act Art.\,52 transparency labelling in their Python source files?

\emph{Hypothesis H1:} Spontaneous Art.\,52 compliance will be below 5\%
of Python files in Copilot-configured repositories, because the Act
provides no enforcement mechanism, no specified format, and no incentive
for early adoption absent regulatory pressure.

\textbf{RQ2 (Label Durability):} Does an Art.\,52 compliance header,
once written, persist through common Python code transformation
operations---formatters, linters, and compilation---that occur in
standard CI/CD pipelines?

\emph{Hypothesis H2:} Standard comment-preserving formatting tools
(\texttt{black}, \texttt{isort}, \texttt{ruff}, \texttt{autopep8},
\texttt{pyupgrade}) will not strip Art.\,52 headers, but compilation to
bytecode and any comment-removal step will destroy them, exposing a
structural durability gap.

\textbf{RQ3 (Supply Chain Provenance):} Is there any evidence of
Art.\,52 compliance markers in the distributed Python software supply
chain---specifically in PyPI package distributions that constitute the
downstream deployment layer of most Python software?

\emph{Hypothesis H3:} Art.\,52 compliance markers will be absent from
PyPI package distributions, demonstrating complete provenance loss at
the distribution layer regardless of what happens at the source level.

% =========================================================================
\section{Experimental Methodology}
\label{sec:methodology}
% =========================================================================

All three experiments use exclusively publicly available, freely
accessible data and introduce no personal data processing.  GitHub data
was accessed via the authenticated GitHub API (public repositories only).
PyPI data was accessed via the public PyPI JSON API and public
distribution downloads.  No user tracking, scraping of private data, or
processing of personal information was performed.

\subsection{Experiment~1---GitHub Mining Study (E1)}

\textbf{Protocol.}  We used the GitHub Code Search API to measure:

(a)~\emph{Total denominator:} The number of public repositories
containing \texttt{.github/copilot-instructions.md}, a file created
exclusively by GitHub Copilot's workspace configuration feature.  This
provides a lower-bound estimate of repositories actively configured for
Copilot use.

(b)~\emph{Compliance signals:} Global counts of public Python files
matching each of five labelling patterns: the exact Art.\,52 citation
(``EU AI Act Art.\,52''), the generic ``AI-Generated Content'' label,
``Generated by:~copilot'', comments referencing GitHub Copilot, and the
\texttt{SPDX-AI-Generated} field proposed for future SPDX versions.

(c)~\emph{Per-repository audit:} A stratified sample of 20~repositories
drawn from the copilot-instructions.md population.  For each, we
queried: total Python file count and count of Python files matching any
Art.\,52 pattern.

Rate limits were respected via enforced inter-request delays.  All
queries targeted public repositories.  No content from private
repositories was accessed.

\textbf{Limitations.}  GitHub code search returns \texttt{total\_count}
estimates that may include false positives (e.g., documentation files
discussing the regulation).  The per-repository sample of $n{=}20$ is
small due to API rate limits; we present it as indicative rather than
definitive and bound our claims accordingly.  The
copilot-instructions.md signal undercounts actual Copilot use (many
teams use Copilot without workspace configuration), making our compliance
rate estimate a ceiling, not a floor.

\subsection{Experiment~2---Header Survival Analysis (E2)}

\textbf{Protocol.}  A canonical Art.\,52-compliant Python file was
generated using pr-automation-agent v0.1.2 (the tool's scaffold output,
selected because it contains a representative five-element Art.\,52
header).  The file was subjected to eleven transformation operations,
each applied independently to the original file:

\begin{enumerate}
\item \texttt{black} (PEP\,8 formatter)
\item \texttt{isort} (import sorter)
\item \texttt{black} + \texttt{isort} (combined)
\item \texttt{ruff format} (Rust-based formatter)
\item \texttt{ruff format} + \texttt{ruff check --fix}
\item \texttt{autopep8 --in-place} (legacy PEP\,8 fixer)
\item \texttt{pyupgrade --py310-plus} (syntax moderniser)
\item \texttt{pyupgrade} + \texttt{black}
\item Full pipeline: \texttt{pyupgrade}$\to$\texttt{isort}$\to$\texttt{black}$\to$\texttt{ruff check --fix}
\item All-comment removal (simulating CI comment-stripping steps)
\item \texttt{python3 -m py\_compile} (compilation to \texttt{.pyc} bytecode)
\end{enumerate}

After each transformation, the output was checked for the three most
critical header elements: the Art.\,52 citation label, the ``Human
reviewer required'' notice, and the ``Generated by:'' attribution field.
Survival was recorded as binary (all three present / not all three present).

\subsection{Experiment~3---PyPI Supply Chain Provenance Loss (E3)}

\textbf{Protocol.}  We selected 25~packages spanning five categories
relevant to the AI-assisted development ecosystem: scientific computing
(\texttt{numpy}, \texttt{pandas}, \texttt{scipy}, \texttt{scikit-learn},
\texttt{matplotlib}); web/API frameworks (\texttt{requests},
\texttt{fastapi}, \texttt{flask}, \texttt{httpx}, \texttt{aiohttp});
developer tooling (\texttt{black}, \texttt{ruff}, \texttt{pytest},
\texttt{mypy}, \texttt{rich}); data engineering---the primary target
domain of AI-generated ingest code (\texttt{dagster}, \texttt{prefect},
\texttt{apache-airflow}, \texttt{dbt-core},
\texttt{great-expectations}); and AI/LLM tooling (\texttt{langchain},
\texttt{openai}, \texttt{anthropic}, \texttt{transformers},
\texttt{datasets}).

For each package, we downloaded the latest source distribution
(\texttt{.tar.gz} or \texttt{.whl}) via the PyPI public JSON API,
extracted all Python files, and applied a regex search for four AI
disclosure patterns: the Art.\,52 citation, the ``AI-Generated Content''
label, any reference to GitHub Copilot in code comments, and the
\texttt{SPDX-AI-Generated} field.  We report per-file and per-package
disclosure rates.

% =========================================================================
\section{Results}
\label{sec:results}
% =========================================================================

\subsection{E1---GitHub Mining: Near-Zero Voluntary Compliance}

Table~\ref{tab:global} reports global Code Search counts;
Table~\ref{tab:sample} reports the per-repository sample.

\begin{table}[htbp]
\caption{GitHub-Scale Art.\,52 Compliance Signals}
\label{tab:global}
\begin{center}
\begin{tabular}{lc}
\toprule
Signal & Global Count \\
\midrule
Repos with \texttt{.github/copilot-instructions.md} & 43,816 \\
Python files: ``EU AI Act Art.\,52''               & 49,248 \\
Python files: ``AI-Generated Content''             & 35,984 \\
Python files: ``Human reviewer required''          & 21,280 \\
Python files: ``Generated by: copilot''            &  2,120 \\
Python files: \texttt{SPDX-AI-Generated}           &      0 \\
\bottomrule
\end{tabular}
\end{center}
\end{table}

\begin{table}[htbp]
\caption{Per-Repository Sample ($n = 20$ Copilot-Configured Repos)}
\label{tab:sample}
\begin{center}
\begin{tabular}{lc}
\toprule
Metric & Value \\
\midrule
Repos in sample                                    & 20          \\
Repos with any Art.\,52 header in Python files    & 2 (10\%)    \\
Repos with any AI disclosure label                 & 2 (10\%)    \\
Python files checked                               & 47          \\
Python files with Art.\,52 header                  & 5 (10.6\%)  \\
Python files with \texttt{SPDX-AI-Generated}       & 0 (0\%)     \\
\bottomrule
\end{tabular}
\end{center}
\end{table}

\textbf{Interpretation.}  43,816 public repositories are actively
configured for GitHub Copilot use.  The global count of Python files
carrying any Art.\,52 citation (49,248) is superficially comparable to
this number, but critically, that global count spans \emph{all} of
public GitHub and includes documentation files, tutorials, academic
papers discussing the regulation, and tools like pr-automation-agent
itself.  The global Art.\,52 count does not represent compliance in
Copilot-using repositories---it represents all public discussion of the
regulation in Python files.  In our directly sampled Copilot
repositories, only 2 of 20 (10\%) had \emph{any} Art.\,52 marker in
their Python files, and 0 of 20 used the \texttt{SPDX-AI-Generated}
format.  The \texttt{SPDX-AI-Generated} count of zero across all of
public GitHub confirms that the one technically machine-verifiable format
proposed to date has achieved zero adoption.

\textbf{H1 assessment:} Confirmed.  Compliance rates in
Copilot-configured repositories are at or below 10\% at the repository
level and 10.6\% at the file level, with the file-level estimate likely
an overestimate given the small denominator.

\subsection{E2---Header Survival: Durable Against Formatters, Destroyed
by Compilation}

Table~\ref{tab:survival} reports header survival results for all eleven
transformations.

\begin{table}[htbp]
\caption{Art.\,52 Header Survival Across Code Transformations.
         Survival rate: 9/11 (81.8\%).}
\label{tab:survival}
\begin{center}
\begin{tabular}{lcc}
\toprule
Transformation & Version & Survives \\
\midrule
Baseline (generated file)                      & ---        & $\checkmark$ \\
\texttt{black}                                 & 26.5.1     & $\checkmark$ \\
\texttt{isort}                                 & current    & $\checkmark$ \\
\texttt{black} + \texttt{isort}                & ---        & $\checkmark$ \\
\texttt{ruff format}                           & 0.15.20    & $\checkmark$ \\
\texttt{ruff format} + \texttt{ruff check}     & 0.15.20    & $\checkmark$ \\
\texttt{autopep8 --in-place}                   & current    & $\checkmark$ \\
\texttt{pyupgrade --py310-plus}                & current    & $\checkmark$ \\
\texttt{pyupgrade} + \texttt{black}            & ---        & $\checkmark$ \\
Full pipeline (pyupgrade$\to$isort$\to$black$\to$ruff) & --- & $\checkmark$ \\
Comment stripping                              & ---        & $\times$     \\
Python \texttt{.pyc} compilation (CPython 3.12.11) & ---    & $\times$     \\
\bottomrule
\end{tabular}
\end{center}
\end{table}

The nine surviving transformations cover the standard Python CI/CD
formatting pipeline.  None of the major formatting or linting tools
(\texttt{black}, \texttt{isort}, \texttt{ruff}, \texttt{autopep8},
\texttt{pyupgrade}) strip code comments, and the Art.\,52 header is
structurally indistinguishable from any other block comment.  Survival
in this sense is not a designed property of the header---it is
incidental.

The two failure modes are structurally significant.  First, Python
compilation to \texttt{.pyc} bytecode strips all comments at compile
time; this is not a bug but a defined property of the CPython compiler.
Any organisation distributing compiled Python (e.g., via wheel
distributions with pre-compiled \texttt{.pyc} files, or via tools like
Cython, Nuitka, or PyInstaller) destroys all Art.\,52 provenance at
distribution time.  Second, CI pipelines that strip comments (common in
minification workflows for embedded Python, infrastructure-as-code
converters, or AI-based code refactoring tools) also destroy the header.

\textbf{H2 assessment:} Confirmed.  The label survives the standard
formatter pipeline but not compilation or comment removal, exposing a
lifecycle durability gap for any Python deployment that involves
compilation or distribution.

\subsection{E3---PyPI Supply Chain: Zero Provenance at Distribution Layer}

Table~\ref{tab:pypi} reports the PyPI audit results across 25~packages
and 18,210 Python source files.

\begin{table*}[htbp]
\caption{Art.\,52 Compliance in PyPI Package Distributions
         ($n = 25$ packages, 18,210 Python files)}
\label{tab:pypi}
\begin{center}
\begin{tabular}{lccccc}
\toprule
Category & Packages & Python Files & Art.\,52 Files &
  Copilot Mention & Disclosure Rate \\
\midrule
Scientific computing & 5 & 5,552 & 0 & 0 & 0.000\% \\
Web/API frameworks   & 5 & 1,475 & 0 & 0 & 0.000\% \\
Developer tooling    & 5 & 2,054 & 0 & 0 & 0.000\% \\
Data engineering     & 5 & 3,761 & 0 & 0 & 0.000\% \\
AI/LLM tooling       & 5 & 5,368 & 0 & 1$^{*}$ & 0.000\% \\
\midrule
\textbf{Total} & \textbf{25} & \textbf{18,210} & \textbf{0} &
  \textbf{1} & \textbf{0.000\%} \\
\bottomrule
\multicolumn{6}{l}{\footnotesize $^{*}$Single Copilot mention (anthropic SDK
  v0.115.1) in a comment unrelated to AI content labelling.}
\end{tabular}
\end{center}
\end{table*}

Across 18,210 Python files from 25 major packages---including those from
organisations (OpenAI, Anthropic, Hugging Face) most likely to be using
AI-assisted development---the Art.\,52 disclosure rate is exactly
0.000\%.  This is not a measurement artefact: these packages collectively
represent billions of monthly downloads and include the foundational
libraries of the Python AI ecosystem.  If Art.\,52 compliance were
occurring at the source level and surviving into distributed packages,
it would appear here.  It does not.

\textbf{H3 assessment:} Confirmed.  Art.\,52 compliance markers are
absent from the distributed Python software supply chain at all levels
of abstraction.

% =========================================================================
\section{Two Under-Examined Regulatory Blind Spots}
\label{sec:blindspots}
% =========================================================================

The experimental results make two structural regulatory problems
empirically concrete.

\subsection{The Provenance Escape Problem (Supply Chain Blind Spot)}

E3 demonstrates that even if every developer in every Copilot-using
repository scrupulously added an Art.\,52 header to every generated
file, that provenance would be invisible to any downstream consumer of
the resulting library.  The PyPI distribution---the artefact that most
users actually install and run---contains no AI content provenance
information.

This is not an edge case.  The software supply chain is the primary
delivery mechanism for Python code.  When a data engineering team uses
\texttt{langchain}, \texttt{dagster}, or \texttt{transformers} in their
pipeline, they are executing millions of lines of Python for which they
have zero knowledge of AI generation history.  If any of those lines
were AI-generated and contain a correctness or security error, the
Art.\,52 label that may have existed in the original source repository
provides no protection to the downstream user---it was never transmitted.

This problem is structurally analogous to the SBOM problem that the SPDX
standard and US Executive Order~14028 attempted to address for component
identification.  SPDX~3.0 has introduced an AI extension profile;
extending it to require an AI content generation field in package
manifests would propagate AI provenance through the supply chain in a
machine-readable, standard format.  The EU AI Act's current text
contains no equivalent mandate.

The analogy to SBOM is also instructive about timelines: SBOM
requirements were announced in the US in 2021 and have still not
achieved widespread industry adoption in 2026.  A comparable AI
provenance requirement, without stronger enforcement mechanisms than
those currently in Art.\,52, should not be expected to achieve
meaningful adoption in less time.

\subsection{The Output-Context Risk Elevation Problem (Risk
Classification Blind Spot)}

The Act's risk stratification (prohibited $\to$ high-risk $\to$
limited-risk $\to$ minimal-risk) is applied to the AI system at its
point of deployment.  GitHub Copilot and AI code generation tools are
classified as limited-risk under the current Annex~III taxonomy, which
reserves high-risk classification for systems directly involved in
critical infrastructure, employment decisions, law enforcement, education
assessment, and similar high-stakes determinations.

This classification is legally defensible but functionally misleading
for the specific case of AI-assisted code generation targeting data
engineering workflows.  E3's package survey includes \texttt{dagster}
(1,742 Python files), \texttt{prefect} (841 files), \texttt{dbt-core}
(220 files), and \texttt{great-expectations} (958 files)---the
foundational tools of production data pipelines that process financial
transactions, health records, and behavioural data at scale.
AI-generated code in these contexts is not a limited-risk tool generating
inconsequential output: it generates the ingest and transformation logic
that directly drives Annex~III high-risk systems (credit scoring,
insurance pricing, employment analytics).

The Act's risk classification does not propagate downstream.  A
limited-risk AI tool can generate code that, when executed, drives a
high-risk system---and the deployer's Art.\,52 obligations remain at the
limited-risk level regardless.  We term this the \emph{output-context
risk elevation problem}: the regulatory treatment of the generating tool
is decoupled from the risk context of its generated output.

This blind spot has a direct practical implication for Art.\,52
enforcement.  The Commission's guidance on GPAI risk assessment should
introduce a concept of downstream risk inheritance: if an AI code
generation tool is deployed in a workflow whose outputs feed a system
that would itself be classified as high-risk under Annex~III, the
deployer's Art.\,52 obligations---and specifically, the human oversight
requirements---should be elevated proportionately.  The Commission's
risk guidance for GPAI providers under Art.\,55 provides a partial
precedent for this kind of contextual risk adjustment.

% =========================================================================
\section{Doctrinal Analysis: Art.~52 Within the EU Regulatory Package}
\label{sec:doctrinal}
% =========================================================================

\subsection{The Principle-Without-Specification Problem: Art.~52
vs.\ GDPR Art.~30}

The EU's GDPR (EU~2016/679)~\cite{gdpr2016} provides the closest
structural parallel to Art.\,52.  Both establish transparency obligations
about processing activities affecting individuals; both are backed by
fines of comparable scale (Art.\,83 GDPR; Art.\,99 EU AI Act).  The
structural difference lies in operational specification.

GDPR Art.\,5(1)(a) establishes the transparency \emph{principle}.
GDPR Arts.\,13--14 specify the \emph{content} of information obligations
(specific enumerated fields including categories of data, purposes,
retention periods, recipients).  GDPR Art.\,30 mandates a \emph{Record
of Processing Activities} with defined fields, retention requirements,
and mandatory availability to supervisory authorities on request.
Art.\,5 without Arts.\,13--14 and 30 would be unenforceable: the
principle would be unarguable in court without a specification of what
transparency requires.

Art.\,52 of the EU AI Act, in its current form, is structurally
equivalent to GDPR Art.\,5(1)(a) alone.  The principle is stated; the
specification is absent.  No Art.\,13-equivalent specifies what fields
an AI content label must contain.  No Art.\,30-equivalent mandates a
record of AI-generated content that must be maintained and made available
to the AI Office on request.  The Commission AI Office can issue guidance
under Art.\,64, but guidance is not binding; it cannot substitute for
Implementing Act specifications.

The GDPR's enforcement history makes the practical consequence clear:
in the early years of GDPR application (2018--2020), when DPAs were
still interpreting Arts.\,13--14 field requirements, enforcement was
inconsistent and fines were low.  Substantive enforcement became possible
only once field requirements were codified through DPA guidance and EDPB
recommendations.  Art.\,52 is currently in the equivalent
pre-specification state, with no equivalent of the EDPB to issue binding
recommendations and no Implementing Act in sight.

\subsection{The Three-Instrument Package: EUAIA + AILD + rPLD}

The AI Liability Directive~\cite{aild2022} (proposed 2022, under
trilogue negotiation) includes, in Article~3, a disclosure-of-evidence
mechanism with teeth: where a plaintiff establishes that an AI system was
involved in causing damage, and where the defendant \emph{refuses or is
unable} to disclose relevant evidence about the AI system's functioning,
the court may presume the AI's causal contribution.  This presumption
shifts the burden of proof in negligence claims---a significant departure
from standard product liability law.

For Art.\,52 specifically, the AILD creates a liability asymmetry that
current compliance practice does not address.  A developer who merges
AI-generated code without maintaining an audit trail (which Art.\,52
does not require) cannot, in subsequent litigation, produce the evidence
needed to rebut the AILD presumption.  The Art.\,52 label, had it been
written and preserved, would constitute exactly the kind of ``relevant
evidence about the AI system's functioning'' that the AILD's disclosure
mechanism contemplates.  But a comment in a source file---mutable,
unretained, absent from compiled distributions (E2, E3)---is not a
durable evidentiary record.

The revised Product Liability Directive~((EU) 2024/2853)~\cite{rpld2024}
closes the previous exclusion of software from product liability.
AI-generated code that causes property damage or personal injury may now
impose liability on the ``manufacturer''---which the rPLD defines to
include any entity that ``substantially modifies'' a product.  Whether
deploying AI-generated code constitutes substantial modification is an
open legal question that will need judicial resolution.  In the interim,
organisations deploying AI-generated code face an uncertain liability
exposure that Art.\,52's current framework does not help them document
or mitigate.

The coherent reading of the three instruments is: Art.\,52 creates the
transparency obligation; the AILD creates the liability exposure for
failing to fulfil it; the rPLD extends that exposure to downstream
software deployments.  The package is structurally sound.  It fails
operationally because Art.\,52 lacks the specification needed to make it
the evidentiary record that AILD Art.\,3 requires.

\subsection{Comparative Analysis: China's Deep Synthesis Regulations}

China's Provisions on the Administration of Deep Synthesis Internet
Information Services (``Deep Synthesis Regulations,'' effective
January~10, 2023)~\cite{china2023} address AI-generated content with
greater technical prescriptiveness than EU Art.\,52.  Article~14
requires that deep synthesis service providers add ``explicit labels''
identifying the content as AI-generated; Article~17 specifies that
labels must include the service provider's identity and the content
type.  Article~20 mandates that platforms capable of identifying deep
synthesis content implement technical detection capabilities.

The comparison is illuminating in two directions.  First, China's
regulations are \emph{more prescriptive} than Art.\,52---they specify
label content fields and impose platform detection obligations that the
EU Act does not.  Second, they are \emph{not technically specified} in
the sense of providing a machine-readable schema: the Cyberspace
Administration of China (CAC) regulations name required fields but do
not provide an encoding standard or a verification mechanism equivalent
to C2PA's cryptographic manifests~\cite{c2pa21}.

Both regulatory regimes, in other words, fall short of the technical
standard that would make compliance auditable.  The difference is one
of degree: China's regulations reduce implementation ambiguity by
specifying \emph{what} a label must contain; the EU's leave even that
open.  Neither requires the \emph{how} that would enable machine
verification.

The practical implication for the EU's forthcoming technical standards
work: the minimum viable improvement is China-level specificity
(enumerate required fields); the technically sound improvement is
C2PA-level verifiability (cryptographic provenance binding).  The
choice between these approaches is a genuine policy trade-off between
implementation burden and verifiability, and it is the decision that the
AI Act's standardisation process should be explicitly addressing.

% =========================================================================
\section{Policy Recommendations}
\label{sec:recommendations}
% =========================================================================

The empirical findings and doctrinal analysis converge on four concrete
recommendations.

\textbf{R1---Mandate a minimum label schema via Implementing Act.}  The
Commission should adopt, under Art.\,73, an Implementing Act specifying
the minimum required fields for an Art.\,52 label in AI-generated source
code: at minimum, a tool identifier, a model identifier, a generation
timestamp, and a reference to the human reviewer obligation.  This is
the minimum needed to make Art.\,52 compliance distinguishable from
incidental text.  A SPDX AI extension field in package manifests should
be included in this specification to address the supply chain provenance
gap~(E3).

\textbf{R2---Establish an audit retention obligation.}  The Act should
be amended or supplemented under Art.\,73 to require that deployers of
GPAI in production software workflows maintain a persistent,
independently verifiable record of AI-generated artefacts, analogous to
GDPR Art.\,30 Records of Processing Activities.  The audit record should
be retained for a period proportionate to risk classification and made
available to the AI Office upon request.  This record would also
constitute the evidentiary basis for the AILD Art.\,3 disclosure
defence.

\textbf{R3---Introduce output-context risk elevation.}  Commission
guidance under Art.\,55 on GPAI risk assessment should introduce the
concept of output-context risk inheritance: AI code generation tools
deployed in workflows whose outputs feed Annex~III high-risk systems
should be subject to Art.\,52 obligations commensurate with the
downstream risk classification.  The Commission's existing technical
guidance on systemic risk assessment for GPAI providers provides a
methodological precedent for this kind of contextual risk adjustment.

\textbf{R4---Require platform-level enforcement for major GPAI
deployers.}  For GPAI code generation tools above a deployment threshold
(e.g., the 100~million users threshold used in Art.\,51 for systemic
risk), the Act should require platform-level enforcement of Art.\,52
labelling---meaning that the tool itself prevents generation of
unlabelled output, or cryptographically signs output in a
C2PA-compatible format.  This would shift the compliance burden from
individual developers (who, as Perry et al.~\cite{perry2023} show,
cannot be relied upon to add labels under time pressure) to the platform
providers best positioned to implement it uniformly.

% =========================================================================
\section{Threats to Validity}
\label{sec:threats}
% =========================================================================

\textbf{Internal validity.}  The GitHub mining (E1) uses
\texttt{total\_count} estimates from GitHub's search index, which may
include files that mention Art.\,52 in comments discussing the
regulation rather than implementing it.  We mitigate this by requiring
the specific citation pattern and by cross-checking with the
per-repository audit.  The per-repository sample ($n{=}20$) is small;
we present the aggregate finding as indicative rather than precisely
generalisable.

\textbf{External validity.}  The PyPI package sample (E3, $n{=}25$) was
selected from high-download packages, which may differ from smaller or
newer packages more likely to use AI-generated code.  If smaller, newer
packages are more likely to include AI disclosure labels, our 0.000\%
finding would be an overestimate of compliance failure.  We consider this
unlikely given that smaller packages have fewer compliance resources, but
acknowledge the selection bias.

\textbf{Construct validity.}  We operationalise Art.\,52 compliance as
the presence of a structured comment header, which is the dominant
implementation approach observed in the literature and in the
pr-automation-agent tool used to generate the E2 test file.  Alternative
implementations (README disclosures, commit message conventions, CI
metadata) would not be captured by our search patterns.  If significant
Art.\,52 compliance occurs through these alternative channels, our E1
measurements would undercount compliance.  We note, however, that such
alternatives would also fail the durability test (E2) and the
supply-chain propagation test (E3).

% =========================================================================
\section{Conclusion}
\label{sec:conclusion}
% =========================================================================

We set out to measure whether EU AI Act Article~52 transparency
obligations are being met in practice across publicly available Python
repositories.  Three experiments on freely accessible data provide a
consistent answer: they are not.  Across 43,816 Copilot-configured
repositories, voluntary compliance is at or below 10\% at the repository
level.  Across 18,210 Python files from 25 major PyPI
distributions---including packages from the leading AI tool
providers---the Art.\,52 compliance rate is 0.000\%.  Where labels are
written, they survive standard code formatting but are destroyed by
bytecode compilation and comment-stripping steps that are routine in
production deployment pipelines.

These findings are not a critique of individual developers' compliance
effort.  They are a predictable consequence of a regulatory text that
establishes a principle without operational specification, mandates
transparency without a verifiable format, and requires human oversight
without a workflow control that enforces it.  Art.\,52, in its current
form, satisfies the political requirement for an AI transparency rule
without providing the technical infrastructure that would make
transparency real.

The path from symbolic to substantive compliance requires four specific
amendments: a minimum label schema via Implementing Act, an audit
retention obligation analogous to GDPR Art.\,30, output-context risk
elevation guidance for GPAI tools in high-risk workflows, and
platform-level enforcement obligations for major GPAI providers.
China's Deep Synthesis Regulations demonstrate that more prescriptive
AI content labelling rules are implementable; the C2PA standard
demonstrates that cryptographically verifiable provenance is achievable.
The EU AI Act's transparency framework has the right objective.  It
needs the operational specification to pursue it.

% =========================================================================
\section{Ethical and Legal Considerations}
\label{sec:ethics}
% =========================================================================

\textbf{GDPR and data protection.}  This study processed no personal
data within the meaning of GDPR Art.\,4(1).  All measurements reported
in Tables~\ref{tab:global}--\ref{tab:pypi} are aggregate statistics
derived from public repository metadata and public package distributions.
No individual user profiles, email addresses, commit author identities,
or other personal identifiers were collected, stored, or reported.  The
per-repository sample (Section~\ref{sec:methodology}) used individual
repository names solely as query keys during data collection; no
repository names are reported in the paper.  GitHub Code Search API
results are aggregate \texttt{total\_count} statistics, which are not
personal data.  EU Regulation 2016/679 does not apply to aggregate,
non-identifying statistics derived from publicly available content.

\textbf{GitHub API Terms of Service.}  All GitHub data was retrieved via
the authenticated GitHub REST API v3, targeting exclusively public
repositories.  Requests were authenticated and rate-limited per GitHub's
API documentation.  GitHub's Terms of Service (Section H.1) explicitly
permits public API use for research purposes subject to rate-limit
compliance, which was maintained throughout.  No web-interface scraping
or bypass of access controls was employed.

\textbf{PyPI data.}  PyPI package distributions were retrieved via the
public PyPI JSON API and public download infrastructure, provided for
automated tooling and research use under the PyPI Terms of Use.
Downloaded distributions were analysed in memory for pattern matching
and immediately discarded; no package source code was retained.

\textbf{Copyright.}  Pattern-matching analysis of source code for
empirical research constitutes analysis, not reproduction, under InfoSoc
Directive Art.\,5(3)(a) (scientific research exception) and analogous
fair use provisions in other jurisdictions.  No substantial reproduction
of any copyrighted source file occurred.

\textbf{Ethical review.}  This study involved no human subjects, no
personal data, and no interaction with individuals.  Consistent with the
ethics guidelines of MSR, ICSE, and FSE for studies using exclusively
public repository metadata and package distributions, the study is
classified as exempt from formal ethics review.

% =========================================================================
\begin{thebibliography}{00}

\bibitem{github2023}
GitHub, ``The State of the Octoverse 2023: AI and Developer
Productivity,'' GitHub, Inc., San Francisco, CA, USA, Tech.\ Rep.,
Oct.\ 2023. [Online]. Available: \url{https://octoverse.github.com/}

\bibitem{perry2023}
N.~Perry, M.~Srivastava, D.~Kumar, and D.~Boneh, ``Do users write more
insecure code with AI assistants?'' in \emph{Proc.\ ACM SIGSAC Conf.\
Computer and Communications Security (CCS)}, Copenhagen, Denmark,
Nov.\ 2023, pp.\ 2785--2799.

\bibitem{edelman1992}
L.~B. Edelman, ``Legal ambiguity and symbolic structures: Organizational
mediation of civil rights law,'' \emph{American Journal of Sociology},
vol.\ 97, no.\ 6, pp.\ 1531--1576, May 1992.

\bibitem{machuletz2020}
A.~Machuletz and R.~B\"{o}hme, ``Multiple purposes, multiple problems:
A user study of consent dialogs after GDPR,'' \emph{Proceedings on
Privacy Enhancing Technologies}, vol.\ 2020, no.\ 2,
pp.\ 481--498, 2020.

\bibitem{chen2021}
M.~Chen et al., ``Evaluating large language models trained on code,''
arXiv preprint arXiv:2107.03374, Jul.\ 2021.

\bibitem{spdx23}
Linux Foundation, ``SPDX Specification v2.3,'' Software Package Data
Exchange (SPDX), Jun.\ 2022. [Online]. Available:
\url{https://spdx.github.io/spdx-spec/v2.3/}

\bibitem{c2pa21}
Coalition for Content Provenance and Authenticity (C2PA), ``C2PA
Technical Specification v2.1,'' Jan.\ 2024. [Online]. Available:
\url{https://c2pa.org/specifications/}

\bibitem{aild2022}
European Commission, ``Proposal for a Directive on adapting
non-contractual civil liability rules to artificial intelligence (AI
Liability Directive),'' COM(2022) 496 final, Sep.\ 2022.

\bibitem{rpld2024}
European Parliament and of the Council, ``Directive (EU) 2024/2853 on
liability for defective products,'' \emph{Official Journal of the
European Union}, L, Nov.\ 2024.

\bibitem{gdpr2016}
European Parliament and of the Council, ``Regulation (EU) 2016/679 on
the protection of natural persons with regard to the processing of
personal data (GDPR),'' \emph{Official Journal of the European Union},
L~119, pp.\ 1--88, May 2016.

\bibitem{china2023}
Cyberspace Administration of China, ``Provisions on the Administration
of Deep Synthesis Internet Information Services,'' effective
Jan.\ 10, 2023.

\bibitem{kirchenbauer2023}
J.~Kirchenbauer, J.~Geiping, Y.~Wen, J.~Katz, I.~Miers, and
T.~Goldstein, ``A watermark for large language models,'' in \emph{Proc.\
Int.\ Conf.\ Machine Learning (ICML)}, Honolulu, HI, USA,
Jul.\ 2023, pp.\ 17061--17084.

\bibitem{veale2021}
M.~Veale and F.~Z. Borgesius, ``Demystifying the Draft EU Artificial
Intelligence Act,'' \emph{Computer Law Review International}, vol.\ 22,
no.\ 4, pp.\ 97--112, 2021.

\bibitem{w3c2024}
W3C, ``Verifiable Credentials Data Model v2.0,'' W3C Recommendation,
Feb.\ 2024. [Online]. Available:
\url{https://www.w3.org/TR/vc-data-model-2.0/}

\bibitem{dora2022}
European Parliament and of the Council, ``Regulation (EU) 2022/2554 on
digital operational resilience for the financial sector (DORA),''
\emph{Official Journal of the European Union}, L~333, pp.\ 1--79,
Dec.\ 2022.

\end{thebibliography}

\end{document}
